\scanpage{2}
déc. 81--janv. 82

\uline{Espaces stratifiés et voisinages côniques}

(ou : \uline{déploiement} \uline{des espaces stratifiés})

\diagram{\begin{tikzcd}[row sep=large,column sep=large]
& \Sigma_{ik} \arrow[rr,hook] \arrow[dl,"(F_{ik})"'] && \Sigma_k & \\
\Sigma_i && \Sigma_{ijk} \arrow[ul,hook] \arrow[rr,hook] \arrow[dl] \arrow[ur,phantom,"\mathrm{cart.}"] \arrow[dr,phantom,"\mathrm{cart.}"] && \Sigma_{jk} \arrow[ul,hook] \arrow[dl,"(F_{jk})"] \\
& \Sigma_{ij} \arrow[ul,"(F_{ij})"] \arrow[rr,hook] && \Sigma_j &
\end{tikzcd}}

$(i<j<k)$
\[
\left\{\begin{array}{ll}
\Sigma_i\ \Sigma_j\ \Sigma_k &: \text{strates déployées}\\
\Sigma_{ij}\ \Sigma_{jk}\ \Sigma_{ik} &: \text{tubes de raccord}\\
&\text{des strates 2 à 2}\\
\Sigma_{ijk} &: \text{jonction en }i,j,k
\end{array}\right.
\]
fibres $F_{ij},F_{jk},F_{ik}$ : fibres de raccord
 des strates 2 à 2.

les flèches $\hookrightarrow$ sont bordantes

les flèches $\longrightarrow$ sont des fibrations \inserted{(propres)}
(cas \unsure{équisingulier})

($F_{ijk}$, fibré sur $F_{ij}$, fibre $F_{jk}$ : \unsure{fibro-jonction}
\unsure{en} $i,j,k$

$\Sigma_i$ est la strate \inserted{déployée} ``âme'' du tube de raccord $\unsure{\Sigma_{ij}}$

$\Sigma_j$ en est la strate déployée ambiante

Cas non singulier : a) $\Sigma_i$ variété à bord, dont le bord est
la réunion des $\Sigma_{ji}$ $(j<i)$

b) $\Sigma_{jk}\longrightarrow\Sigma_k$ fibration : fibres des
variétés à bord \inserted{compactes}, le fibré des bords est
la réunion des $\Sigma_{ijk}$ pour $i<j<k$.

Plus généralement : a) $\Sigma_{i_0i_1\ldots i_n}=\Sigma_d$ est une
variété à bord, dont le bord est la réunion des
$\Sigma_{d'}$, avec $d'$ drapeau $\supset d$, cofinal à $d$
(il suffit de les prendre de la forme $d'=d\cup\{j\}$,
avec $j\in I_{<i_0}\amalg I_{i_0<\cdot<i_1}\amalg\cdots\amalg\unsure{I_{i_{n-1}<\cdot<i_n}}$)

b) $\Sigma_{i_0\ldots i_n}=\Sigma_d\longrightarrow\Sigma_{i_0\ldots i_p}=\Sigma_\delta$
est \unsure{fibrant|fibré}, à fibres des variétés à bord compactes,
le fibré des bords est réunion des $\Sigma_{d'}$
avec $d'\supset d$, $d$ cofinal dans $d'$, $\delta$ segment \unsure{initial}
de $d'$ -- il suffit de prendre $d'=d\cup\{j\}$, avec
$j\in I_{i_p<\cdot<\unsure{i_{p+1}}}\amalg\cdots\amalg I_{i_{n-1}<\cdot<i_n}$.

\mspage{1}
